% 第三问圆心测向：构造扇形覆盖圆，并绘制保守半径上界。
% 左图半角放大至18度，整体旋转20度；右图计算仍用1.005度。
% 示意角度与理论数据的说明放在正文图注中。
\newcommand{\QThreeContractionFigure}{\resizebox{\linewidth}{!}{%
\begin{tikzpicture}[>=Stealth,line cap=round,line join=round,
  font=\fontsize{10}{12.5}\selectfont,text=qoneink,
  lab/.style={fill=none,inner sep=2pt,align=center},
  lead/.style={line width=.65pt},
  point/.style={circle,inner sep=1.7pt,fill=qoneink}]
\path[use as bounding box] (0,0) rectangle (16.6,8.0);
\fill[white] (0,0) rectangle (16.6,8.0);
\draw[black!12] (7.9,1.0)--(7.9,7.30);
\node[font=\bfseries] at (3.85,7.6) {(a) 在圆心再次测量，缩小定位区域};
\node[font=\bfseries] at (12.15,7.6) {(b) 多次测量后的半径上界};

% ---------- 左：原包围圆、误差扇形及扇形覆盖圆 ----------
\pgfmathsetmacro{\rad}{2.3}
\pgfmathsetmacro{\coverR}{\rad/(2*cos(18))}
\coordinate (z) at (3.00,4.15);
\coordinate (c) at ($(z)+(20:\coverR)$);
\coordinate (a) at ($(z)+(2:\rad)$);
\coordinate (b) at ($(z)+(38:\rad)$);

\fill[qoneblue!5] (z) circle[radius=\rad];
\draw[qoneblue!85,dash pattern=on 3.5pt off 2.5pt,line width=.85pt]
  (z) circle[radius=\rad];
\fill[qonepurple!7] (c) circle[radius=\coverR];
\fill[qoneorange!27] (z)--(a)
  arc[start angle=2,end angle=38,radius=\rad]--cycle;
\draw[qonepurple,line width=1.0pt] (c) circle[radius=\coverR];
\draw[qoneorange,line width=1.05pt] (a)--(z)--(b);
\draw[qoneorange,line width=1.05pt] (a)
  arc[start angle=2,end angle=38,radius=\rad];

\draw[qoneorange,line width=.65pt] ($(z)+(2:.55)$)
  arc[start angle=2,end angle=38,radius=.55];
\node[text=qoneorange,inner sep=1pt] at ($(z)+(.89,.31)$) {$2\bar\varepsilon$};
\draw[->,qoneblue!85,line width=.7pt] (z)--($(z)+(140:\rad)$);
\node[lab,text=qoneblue] at ($(z)+(-1.12,.54)$) {$r$};
\node[point] at (z) {};
\node[lab,anchor=north east] at ($(z)+(-.10,-.16)$) {检测点 $z$};
\foreach \p in {a,b} \fill[qoneorange] (\p) circle[radius=.042];

\node[lab,text=qoneblue] at (2.10,6.85) {原包围圆 $B(z,r)$};
\node[lab,text=qoneorange,anchor=west,align=left] (sectorlabel) at (5.70,6.20)
  {新区域位于\\此扇形内};
\draw[lead,qoneorange!85] (sectorlabel.south west)--(5.70,5.68)--(5.18,4.93);
\node[lab,text=qonepurple,anchor=west] (coverlabel) at (5.28,2.55) {能罩住扇形的圆};
\draw[lead,qonepurple] (coverlabel.west)--(5.02,2.55)--($(c)+(-49:\coverR)$);
\node[text=qonepurple,font=\fontsize{11}{14}\selectfont] at (3.85,1.13)
  {$\displaystyle r_{\rm new}\le R=\frac{r}{2\cos\bar\varepsilon}$};

% ---------- 右：线性坐标，递推数据未作平滑或修改 ----------
\coordinate (origin) at (8.95,1.55);
\foreach \v in {40,80,120,160} {
  \draw[black!12,line width=.4pt] (8.95,{1.55+\v*.029})--(16.05,{1.55+\v*.029});
}
\draw[->,qoneink!65,line width=.7pt] (origin)--(16.20,1.55);
\draw[->,qoneink!65,line width=.7pt] (origin)--(8.95,6.85);
\node[anchor=west] at (9.02,6.94) {半径上界 / m};
\node at (12.55,.42) {圆心处的检测次数 $k$};
\foreach \v in {0,40,80,120,160} {
  \node[anchor=east,font=\fontsize{9.5}{12}\selectfont]
    at (8.80,{1.55+\v*.029}) {\v};
}

\pgfmathsetmacro{\thresholdY}{1.55+19.5*.029}
\draw[qonefail,dash pattern=on 3pt off 2pt,line width=.8pt]
  (8.95,\thresholdY)--(16.05,\thresholdY);
\node[lab,anchor=south west,text=qonefail,font=\fontsize{9}{11}\selectfont]
  at (9.15,{\thresholdY+.07}) {清除所需半径：$19.5$ m};

\foreach \k in {0,...,4} {
  \pgfmathsetmacro{\bound}{160/(2*cos(1.005))^\k}
  \coordinate (p\k) at ({9.15+\k*1.57},{1.55+\bound*.029});
  \draw[qoneink!50,line width=.5pt] ({9.15+\k*1.57},1.55)--({9.15+\k*1.57},1.44);
  \node[anchor=north,font=\fontsize{9.5}{12}\selectfont]
    at ({9.15+\k*1.57},1.31) {\k};
}
\draw[qoneblue,line width=1.0pt] (p0)--(p1)--(p2)--(p3)--(p4);
\foreach \k in {0,1,2,3}
  \node[point,fill=qoneblue,draw=white,line width=.5pt,inner sep=2.2pt] at (p\k) {};
\node[point,fill=qonecover,draw=white,line width=.5pt,inner sep=2.5pt] at (p4) {};

\node[lab,anchor=south west,text=qoneblue] at ($(p0)+(.02,.05)$) {$160.00$};
\node[lab,anchor=south west,text=qoneblue] at ($(p1)+(.03,.06)$) {$80.01$};
\node[lab,anchor=south west,text=qoneblue] at ($(p2)+(.03,.08)$) {$40.01$};
\node[lab,anchor=south,text=qoneblue] at (13.86,3.10) {$20.01$};
% 末端两点靠近门限，数值放在曲线上方，用短引线对应，不占横轴空间。
\draw[lead,qoneblue!70] ($(p3)+(0,.04)$)--(13.86,3.09);
\node[lab,anchor=south,text=qonecover] at (15.46,3.10) {$10.01$};
\draw[lead,qonecover!80] ($(p4)+(0,.05)$)--(15.43,3.09);
\end{tikzpicture}%
}}
